%%%%%%%%%%%%%%%%%%%%%%%%%%%%%%%
%% Map Areas
%%%%%%%%%%%%%%%%%%%%%%%%%%%%%%%
\ExplSyntaxOn

\RpgThemeKeys{ rpg / map }
{
	\RpgThemeTLKey	{header-offset}	{\__rpg_mapHeaderOffset}{0}
	\RpgThemeTLKey	{title}			{\__rpg_mapTitle}		{}
	\RpgThemeTLKey	{prefix}		{\__rpg_mapPrefix}		{Area~}
	\RpgThemeTLKey	{empty-prefix}	{\__rpg_mapEmptyPrefix}	{Area~}
}
\RpgThemeKeys{ rpg / area }
{
	\RpgThemeTLKey	{insert}		{\__rpg_areaInsertText}	{}
	\RpgThemeTLKey	{label}			{\__rpg_areaLabel}		{}
}

%% Activates an environment designed for describing map areas
%% Complexity comes from the nesting behaviour for names 
\NewDocumentEnvironment{RpgMap}{ O{} }
{
	\group_begin:
	\stepcounter{RpgAreaDepth}			%% Increment nesting depth
	\keys_set:nn{ rpg / map }{#1}
	\int_compare:nT {\value{RpgAreaDepth} = 1 }
	{
		\tl_if_empty:eF{\__rpg_mapTitle}
		{
			\__rpg_mapRenderHeadings:nn{\__rpg_mapHeaderOffset}{\__rpg_mapTitle}
		}
	}
	\__rpg_mapMakePrefix:{}

	%%%%%%%%%%%%
	% \area
	%%%%%%%%%%%%
	% Inside a map, can use \area like you would \item in a list.
	\DeclareDocumentCommand{\area}{O{} g}{\__RpgArea{##1}{##2}}
}
{	\seq_gpop_right:NN \__rpg_MapStack {\l__map_tmp} %%remove the final element from the stack
	\addtocounter{RpgAreaDepth}{-1}		%% Decrement nesting depth

	\group_end:
}

%% Access refs (a wrapper for \ref/\nameref)
%% #1 If starred, call ref, else call nameref
%% #2 The map reference to access
\NewDocumentCommand\RpgMapRef{s m}
{
	\IfBooleanTF{#1}
	{
		\ref{#2}
	}
	{
		\nameref{#2}
	}
}

%% Theme setter function for setting the name formatting system (i.e. 1A-ii)
%% Args: #1 the current nesting depth
%%		 #2 the number to format
\NewDocumentCommand\RpgThemeMapFormat{ +m }
{
	\cs_set:Nn \__rpg_mapFormatAtDepth:nn {#1}
}

%% A nesting environment, equivalent to calling \area[]{} \begin{RpgMap} \end{RpgMap}
%%	#1	Label name for the top level area
%%	#2	Display name for the top level area
\NewDocumentEnvironment{RpgArea}{O{} g}
{
	\int_compare:nT{ \value{RpgAreaDepth} = 0}
	{
		\msg_error:nn{rpg}{bad-map}
	}
	\area[#1]{#2}
	\begin{RpgMap}
		}{
	\end{RpgMap}
}

%%%%%%%%%%%%%%%%%%%%
%% Helper Functions
%%%%%%%%%%%%%%%%%%%%

\newcounter{RpgAreaDepth} %Tracks how many nestings we have
\newcounter{RpgArea}	  %Tracks area references
\setcounter{RpgAreaDepth}{0}
\setcounter{RpgArea}{0}

\seq_new:N \__rpg_MapStack %% The length of this array is the nesting depth, the value at each element is the current no. of elements at that depth
\msg_new:nnn{rpg}{bad-map}{Cannot~activate~an~area~without~an~RpgMap~environment}

% \begin{noindent}
%% Dispatch function for \area and RpgArea
%%	#1	Label name for the area
%%	#2	Display name for the area
\NewDocumentCommand\__RpgArea{m m}
{
	\__rpg_mapIncrementStack:
	\group_begin:
	
	\__rpg_mapMakeHeader:n{#2}
	\__rpg_mapMakeLabels:nn{#1}{\__rpg_labelName}
	\group_end:
}
% \end{noindent}

%% Each area constructs a parent prefix used for descending 
\cs_new_protected:Npn \__rpg_mapMakePrefix:
{
	%% If level = 1, then we're at the top level of a new map, so clear the stack
	\int_compare:nT { \value{RpgAreaDepth} = 1 }
	{
		\seq_gclear:N \__rpg_MapStack
	}

	%% Construct the parent name by iterating over the current stack
	\tl_set:Nn \__rpg_mapParentName  {}
	\seq_map_indexed_inline:Nn \__rpg_MapStack
	{
		\tl_put_right:Ne \__rpg_mapParentName {\__rpg_mapFormatAtDepth:nn{##1}{##2}}
	}

	%% Then increment the stack with the new element
	\seq_gput_right:Nn \__rpg_MapStack {0} %% no elements at this depth, so counter starts at 0
}


%% Using \RpgThemeMapFormat to set the default 
\cs_new_protected:Npn \__rpg_mapFormatAtDepth:nn {}
\RpgThemeMapFormat{
	\int_case:nnF{ \int_mod:nn{#1 - 1} { 3 } }
	{
		{0}{	\int_to_arabic:n{#2}	}
			{1}{	\int_to_Alph:n{#2}	}
			{2}{	-\int_to_roman:n{#2}	}
	}
	{}
}

%Extract the current-depth counter, increment it, and then insert the new value
%(Weirdly is the most optimal way to increment a value in a sequence)
\cs_new_protected:Npn \__rpg_mapIncrementStack:
{
	\seq_gpop_right:NN \__rpg_MapStack \l__tmpa_tl
	\tl_set:Ne \__rpg_MapCurrentDepth { \int_eval:n {1 + \l__tmpa_tl} }
	\seq_gput_right:Ne \__rpg_MapStack {\__rpg_MapCurrentDepth}
}

% Creates the area heading (or sets a default name)
\cs_new_protected:Npn \__rpg_mapMakeHeader:n#1
{
	\tl_set:Ne{\__rpg_mapElementID}{\__rpg_mapParentName\__rpg_mapFormatAtDepth:nn{\arabic{RpgAreaDepth}}{\__rpg_MapCurrentDepth}}

	\IfValueTF{#1}
	{
		\tl_set:Ne{\__rpg_mapElementName}{\__rpg_mapPrefix \__rpg_mapElementID :~#1}
		\def\__rpg_labelName{#1}
	}
	{
		\tl_set:Ne{\__rpg_mapElementName}{\__rpg_mapEmptyPrefix \__rpg_mapElementID}
		\def\__rpg_labelName{\__rpg_mapEmptyPrefix \__rpg_mapElementID}
	}

	%% Actually write the headings -- getting smaller as the nesting increases
	\__rpg_mapRenderHeadings:nn{\arabic{RpgAreaDepth} + \__rpg_mapHeaderOffset}{\__rpg_mapElementName}
}


%% Choose which \section/\subsection to write the header with
% \begin{noindent}
\cs_new_protected:Npn \__rpg_mapRenderHeadings:nn#1#2
{
	\def\__rpg_inserter{}
	\str_case_e:nnF{\int_eval:n{#1}}
	{
		{0}{\let\rpgtmp\chapter}
		{1}{\let\rpgtmp\section}
		{2}{\let\rpgtmp\subsection}
		{3}{\let\rpgtmp\subsubsection}
		{4}{\let\rpgtmp\paragraph \def\__rpg_inserter{\mbox{}}}
	}
	{
		\let\rpgtmp\paragraph \def\__rpg_inserter{\mbox{}}
	}

	\rpgtmp*{#2}\__rpg_inserter
}

%% Generate the labels so the \ref and \nameref work as expected
\cs_new_protected:Npn \__rpg_mapMakeLabels:nn#1#2
{
	\tl_if_empty:eF{#1}
	{
		\refstepcounter{RpgArea}
		\tl_set:Ne \@currentlabelname{#2}
		\tl_set:Ne\@currentlabel{\__rpg_mapElementID}
		\label{#1}	 
	}
}
% \end{noindent}
